% Generated draft: requires mathematical and source review.
\documentclass[11pt]{article}
\usepackage[margin=25mm]{geometry}
\usepackage{fontspec}
\setmainfont{Latin Modern Roman}
\usepackage{amsmath,amssymb,mathtools,unicode-math}
\setmathfont{Latin Modern Math}
\usepackage{enumitem,needspace,xurl,hyperref}
\setlength{\parindent}{0pt}
\setlength{\parskip}{5pt}
\setlength{\emergencystretch}{3em}
\newcommand{\N}{\mathbb{N}}
\newcommand{\Z}{\mathbb{Z}}
\newcommand{\Q}{\mathbb{Q}}
\newcommand{\R}{\mathbb{R}}
\newcommand{\C}{\mathbb{C}}
\newcommand{\F}{\mathbb{F}}
\newcommand{\A}{\mathbb{A}}
\newcommand{\PP}{\mathbb{P}}
\newcommand{\E}{\mathbb{E}}
\newcommand{\eps}{\varepsilon}
\newcommand{\dd}{\,\mathrm{d}}
\begin{document}
\section*{UnsolvedMath ID: 30005640}
\textbf{Source problem code:} catalogue entry
\subsection{Definitions and mathematical statement}
Let \(k\) be a number field and let \(C\) be a smooth curve over \(k\) with function field \(K=k(C)\). Here integrality of \(C\) is implicit in the use of its function field; the supplied record does not explicitly require properness or geometric integrality.

A \(K\)-torus is an algebraic group \(T\) over \(K\) such that \(T_{K^{\mathrm{sep}}}\cong(\mathbb G_m)^r\) for some integer \(r\geq 0\). A group of multiplicative type is a commutative affine group scheme that becomes diagonalizable over a separable closure. The supplied generalization does not specify whether such groups must be of finite type.

For the usual interpretation using closed points and completions, let \(C^{(1)}\) be the set of closed points of \(C\), and let \(K_c\) be the completion of \(K\) at the discrete valuation associated with \(c\in C^{(1)}\). Define
\[
\operatorname{Sha}_C^1(T)
:=\ker\!\left(H^1_{\mathrm{fppf}}(K,T)
\longrightarrow\prod_{c\in C^{(1)}}H^1_{\mathrm{fppf}}(K_c,T)\right).
\]
These cohomology groups classify torsors; for a torus they agree with the usual degree-one Galois cohomology groups. The supplied record does not itself define \(\operatorname{Sha}_C^1\), so this local interpretation remains a qualification requiring source confirmation.

\textbf{Finiteness question.} For every number field \(k\), every smooth curve \(C/k\) as above, and every \(K\)-torus \(T\), is
\[
\#\operatorname{Sha}_C^1(T)<\infty?
\]
The supplied record also asks this for every group of multiplicative type over \(K\); the intended scope of that extension, particularly its finite-type hypothesis, requires clarification. The displayed formulation uses exactly the points of the given curve \(C\), without replacing it by a proper model.

\textbf{Supplied qualification.} The original excerpt says that the conjecture is known for a stably rational \(K\)-torus \(T\), referring to an unspecified ``loc. cit.'' Here stable rationality means that \(T\times_K\mathbb A_K^m\) is \(K\)-birational to \(\mathbb A_K^{\dim(T)+m}\) for some integer \(m\geq0\). The supplied assessment also mentions unspecified related cases and labels the problem partially solved. These are supplied literature claims, not independently verified conclusions.
\subsection{Short English statement}
For a torus over the function field of a smooth curve over a number field, ask whether there are only finitely many torsor classes that become trivial over every completion attached to a closed point of the curve. The source also proposes the question for groups of multiplicative type. The local convention and the scope of this broader class require confirmation. The supplied curation reports the stably rational case as known.
\subsection{Sources}
\begin{itemize}
\item Local UnsolvedMath registry: supplied complete source record, permanent ID \(30005640\), title \emph{Tate-Shafarevich Finiteness for Tori over Function Fields}, exact source problem code \texttt{catalogue entry}.
\item Supplied report citation: \emph{MFO--RIMS Tandem Workshop: Arithmetic Homotopy and Galois Theory (2024)}, Oberwolfach Reports, pp.\ 2377--2488, DOI \texttt{10.4171/owr/2023/42}. Supplied source link: \url{https://doi.org/10.4171/owr/2023/42}. The record associates the problem with 2023, while its report title specifies 2024; this discrepancy is unresolved. No individual proposer is supplied.
\item Supplied literature source: ``Harari, finiteness theorems in Galois cohomology,'' \url{https://doi.org/10.4171/owr/2023/42}. This link identifies the same workshop report; the record does not give an individual contribution's page range or resolve the original excerpt's ``loc. cit.'' reference.
\item Literature checked at: \texttt{2026-08-22}. This exact date belongs to the supplied dated curation, not an independently verified assessment. That curation labels the problem partially solved, reports finiteness for stably rational tori and unspecified related cases, and describes arbitrary tori as conjectural. It explicitly requests independent verification of time-sensitive status and cited claims.
\item The registry labels the statement \texttt{corrected\_verified} and supplies a claim that it was checked against the official report. This draft preserves that provenance claim without asserting an independent source visit or mathematical verification.
\end{itemize}
\end{document}
